\documentclass[11pt]{article}
\linespread{1.5}
\usepackage{fullpage}
\usepackage{amsmath,theorem,amssymb,graphicx,pgfplots,tabularx,placeins}
\usepackage[semicolon,authoryear]{natbib}
\usepackage{caption}
\usepackage{subcaption}
\usepackage{csquotes}
\usepackage{epstopdf}
\pgfplotsset{compat=1.18}

\newcommand{\lb}{\label}
\newtheorem{thm}{Theorem}
\newtheorem{prop}{Proposition}
\newtheorem{definition}{Definition}

\newcommand{\bit}{\begin{itemize}}
	\newcommand{\eit}{\end{itemize}}
\newcommand{\ben}{\begin{enumerate}}
	\newcommand{\een}{\end{enumerate}}
\newcommand\setItemnumber[1]{\setcounter{enumi}{\numexpr#1-1\relax}}

\title{Draft Solutions and Marking Guide\\
Ph.D. Macroeconomics II, Questions 2--4}
\author{}
\date{}

\begin{document}
\maketitle

\section*{Question 2: Government spending in the New-Keynesian model}

\ben
\item Goods clearing implies
\[
\hat c_t=\frac{\hat y_t-s_G\hat g_t}{1-s_G}.
\]
Using $\hat n_t=\hat y_t$ and $\widehat{mc}_t=\hat c_t+\varphi \hat n_t$,
\[
\widehat{mc}_t=
\left(\frac{1}{1-s_G}+\varphi\right)\hat y_t
-\frac{s_G}{1-s_G}\hat g_t.
\]
Flexible prices imply $\widehat{mc}_t=0$, hence
\[
\hat y^n_t=
\frac{s_G}{1+\varphi(1-s_G)}\hat g_t.
\]
Give up to 2 points for goods clearing/private consumption, 2 points for marginal cost, and 1 point for natural output.

\item Let
\[
D\equiv 1+\varphi(1-s_G), \qquad \psi\equiv \frac{s_G}{D}.
\]
Then $\hat y_t^n=\psi \hat g_t$. The natural level of private consumption is
\[
\hat c_t^n
=\frac{\hat y_t^n-s_G\hat g_t}{1-s_G}
=-\frac{\varphi s_G}{D}\hat g_t.
\]
The natural real interest rate is
\[
r_t^n=E_t\hat c^n_{t+1}-\hat c^n_t
=(1-\rho_g)\frac{\varphi s_G}{D}\hat g_t.
\]
The output gap is $x_t=\hat y_t-\hat y_t^n$. Since the private-consumption gap is $x_t/(1-s_G)$, the gap IS equation is
\[
x_t=E_t x_{t+1}-(1-s_G)(\hat i_t-E_t\pi_{t+1}-r_t^n).
\]
The Phillips curve is
\[
\pi_t=\beta E_t\pi_{t+1}+\kappa_y x_t,
\qquad
\kappa_y\equiv \kappa\left(\frac{1}{1-s_G}+\varphi\right)
=\kappa\frac{D}{1-s_G}.
\]
Give 2 points for defining the gap and the Phillips curve, 2 points for the correctly scaled IS equation, and 1 point for $r_t^n$.

\item Guess $x_t=a_x\hat g_t$ and $\pi_t=a_\pi\hat g_t$. Since $E_t\hat g_{t+1}=\rho_g\hat g_t$,
\[
(1-\rho_g)a_x+(1-s_G)(\phi_\pi-\rho_g)a_\pi
=(1-s_G)(1-\rho_g)\frac{\varphi s_G}{D},
\]
and
\[
(1-\beta\rho_g)a_\pi=\kappa_y a_x.
\]
Thus
\[
a_x=
\frac{(1-s_G)(1-\rho_g)\frac{\varphi s_G}{D}}
{(1-\rho_g)+(1-s_G)(\phi_\pi-\rho_g)\frac{\kappa_y}{1-\beta\rho_g}},
\]
and
\[
a_\pi=\frac{\kappa_y}{1-\beta\rho_g}a_x.
\]
Both coefficients are positive under the stated assumptions. A higher $\phi_\pi$ raises the denominator and therefore lowers $a_x$ and $a_\pi$. Give 3 points for the algebraic solution and 2 points for signs/comparative statics.

\item Under flexible prices, $x_t=0$, $\pi_t=0$, and output rises by
\[
\hat y^n_t=\frac{s_G}{D}\hat g_t.
\]
Private consumption falls because government purchases absorb resources and create a negative wealth effect:
\[
\hat c_t^n=-\frac{\varphi s_G}{D}\hat g_t<0.
\]
Under sticky prices and the Taylor rule, the government-spending shock raises the natural real rate. Since the nominal-rate rule does not perfectly replicate the natural real rate, the real rate is too low relative to $r_t^n$, so $x_t>0$ and $\pi_t>0$. Output rises by more than under flexible prices:
\[
\hat y_t=\hat y_t^n+x_t.
\]
Private consumption may fall by less than under flexible prices, and can in principle rise if the positive output gap is large enough, since
\[
\hat c_t=\hat c_t^n+\frac{x_t}{1-s_G}.
\]
Optimal monetary policy stabilizes the output gap and inflation, implementing $x_t=0$ and $\pi_t=0$ by setting the real rate equal to the natural real rate. Give 2 points for flexible prices, 2 points for sticky-price mechanism, and 1 point for optimal policy.
\een


\section*{Question 3: Reservation wages and wage dispersion}

\ben
\item The Bellman equations are
\[
rW(w)=w+\sigma[U-W(w)],
\]
\[
rU=b+\lambda_u\int \max\{W(w)-U,0\}dF(w).
\]
Since $W(w)$ is increasing in $w$, there is a cutoff $w_R$ such that the worker accepts iff $w\geq w_R$, with $W(w_R)=U$. Give 2 points for Bellman equations and 2 points for cutoff logic.

\item From the employed Bellman equation,
\[
W(w)=\frac{w+\sigma U}{r+\sigma}.
\]
Using $W(w_R)=U$ gives $rU=w_R$. Hence
\[
W(w)-U=\frac{w-w_R}{r+\sigma}.
\]
Substitute into the unemployment Bellman equation:
\[
w_R=b+\lambda_u\int_{w\geq w_R}\frac{w-w_R}{r+\sigma}dF(w),
\]
or
\[
w_R-b=\frac{\lambda_u}{r+\sigma}\int_{w\geq w_R}(w-w_R)dF(w).
\]
The reservation wage increases with $b$ and $\lambda_u$, and decreases with $\sigma$. Intuitively, higher benefits or more frequent offers make waiting more valuable; faster separations make accepted jobs less valuable. Give 3 points for derivation and 2 points for comparative statics.

\item Since
\[
\int_{w\geq w_R}(w-w_R)dF(w)
=Ew-w_R-\int_{\underline w}^{w_R}(w-w_R)dF(w),
\]
and
\[
-\int_{\underline w}^{w_R}(w-w_R)dF(w)
=\int_{\underline w}^{w_R}(w_R-w)dF(w)
=\int_{\underline w}^{w_R}F(w)dw,
\]
we obtain
\[
w_R-b=\frac{\lambda_u}{r+\sigma}
\left[
Ew-w_R+\int_{\underline w}^{w_R}F(w)dw
\right].
\]
For a given reservation wage, a mean-preserving spread increases the option value of search because the worker can reject sufficiently bad offers while preserving the upside from very good offers. Therefore the right-hand side of the reservation-wage equation rises at the old reservation wage, and the equilibrium reservation wage increases. Give 2 points for the transformation and 2 points for the economic interpretation.

\item A higher reservation wage lowers the acceptance probability $1-F(w_R)$ and therefore increases expected unemployment duration. Accepted wages are drawn from a higher part of the offer distribution, so average accepted wages and match quality rise. Give 1 point for unemployment duration and 1 point for accepted wages.
\een


\section*{Question 4: Idiosyncratic risk in the Aiyagari model}

\ben
\item The household problem is
\[
V(a,\epsilon)=\max_{c,a'}\left\{
u(c)+\beta E[V(a',\epsilon')|\epsilon]
\right\}
\]
subject to
\[
c+a'=w\epsilon+(1+r)a,\qquad a'\geq 0,\qquad c\geq 0.
\]
The firm solves
\[
\max_{K,L} K^\alpha L^{1-\alpha}-(r+\delta)K-wL.
\]
The first-order conditions are
\[
r+\delta=\alpha K^{\alpha-1}L^{1-\alpha},
\qquad
w=(1-\alpha)K^\alpha L^{-\alpha}.
\]
With mean labor efficiency equal to one, $L=1$ in stationary equilibrium. Give 2 points for household problem and 2 points for firm FOCs.

\item A stationary recursive competitive equilibrium consists of prices $(r,w)$, policy functions $c(a,\epsilon)$ and $a'(a,\epsilon)$, firm choices $(K,L)$, and an invariant distribution $\mu(a,\epsilon)$ such that: households solve their recursive problem; firms optimize; $\mu$ is invariant under the household policy and the Markov process for $\epsilon$; labor clears with $L=\int \epsilon d\mu=1$; asset/capital markets clear with
\[
K=\int a\,d\mu;
\]
and the goods market clears,
\[
\int c(a,\epsilon)d\mu+\delta K=K^\alpha L^{1-\alpha}.
\]
Give up to 4 points for these equilibrium components.

\item The Aiyagari diagram plots aggregate asset supply $A^s(r)$ from household saving behavior and capital demand $K^d(r)$ from firm optimality. Capital demand is downward sloping:
\[
K^d(r)=\left(\frac{\alpha}{r+\delta}\right)^{1/(1-\alpha)}
\]
when $L=1$. Aggregate asset supply is increasing in $r$: a higher return makes saving more attractive. With uninsured idiosyncratic risk and prudence, households have a precautionary-saving motive. As $r$ approaches $1/\beta-1$ from below, desired saving becomes very large. A stationary equilibrium therefore requires
\[
r<1/\beta-1.
\]
Give 1 point for demand, 1 point for supply/diagram, and 1 point for the interest-rate bound.

\item More idiosyncratic income risk strengthens precautionary saving. For any given $r$, households want to hold more assets, so $A^s(r)$ shifts to the right. In equilibrium, the interest rate falls and the capital stock rises. The higher capital stock raises wages and lowers the marginal product of capital. In a representative-agent or complete-markets economy with the same mean labor endowment, purely idiosyncratic risk is insured or diversified away, so this precautionary-saving channel is absent. Give 2 points for the shift and equilibrium effects, and 2 points for the complete-markets contrast.
\een

\end{document}
